% GENERATED FILE. Edit canonical lessons, then run npm run generate:books.
% canonical-source-hash: fnv1a64:17110a0cf3ffd132
% canonical-lessons: ES-C609-progreso, ES-C609-perfeccionar, ES-C609-desafio, ES-C609-afrontar, ES-C609-perseverar

\chapter{Progresar}
\label{ch:es-progresar}
\hlchaptermodality{609}

\begin{chapteropening}
\textbf{By the end of this chapter:} \emph{I can talk about progress, a challenge and how to face it and do better.}

Complete the last lesson of chapter 609: I can talk about progress, a challenge and how to face it and do better.
\end{chapteropening}

\section[el progreso]{el progreso — the progress}

\label{lesson:ES-C609-progreso}



\begin{quote}
Before the new one: say the Spanish for the growth, then the Spanish for the trend.
\end{quote}

\subsection*{You'll want to know: el progreso}
\textbf{el progreso} — \textquotedblleft{}the progress\textquotedblright{}.

\begin{quote}
Has hecho muchos progresos.
\end{quote}

\begin{quote}
\emph{You have made a lot of progress.}
\end{quote}

\emph{Hacer progresos} — to make progress; the plural is usual here.

\begin{cousinweb}
\emph{progreso} is a learned borrowing from Latin \emph{progressus} 'an advance', built on \emph{progredi} 'to go forward', from \emph{gradi} 'to step' (Lewis \& Short).
\end{cousinweb}

\begin{grammarlens}[title={What you've built}]
One more word for telling, explaining and describing at B1.
\end{grammarlens}

\subsection*{Your turn}
\begin{itemize}
\raggedright
  \item \emph{Say it:} \emph{el progreso}
  \item \emph{Say it:} \emph{el progreso}, once more
  \item \emph{Say it:} a sentence of your own with \emph{el progreso}
  \item \emph{Recall:} say the Spanish for the growth, then the Spanish for the trend, then say \emph{el progreso} again
\end{itemize}

\begin{tcolorbox}[breakable,colback=teal!4,colframe=teal!35!black,title={Before you move on}]
What does el progreso mean? (\textquotedblleft{}The progress\textquotedblright{}.) Say it once more.
\end{tcolorbox}
\section[perfeccionar]{perfeccionar — to perfect, to polish}

\label{lesson:ES-C609-perfeccionar}



\begin{quote}
Before the new one: say the Spanish for the trend, then the Spanish for the progress.
\end{quote}

\subsection*{You'll want to know: perfeccionar}
\textbf{perfeccionar} — \textquotedblleft{}to perfect, to polish\textquotedblright{}.

\begin{quote}
Quiero perfeccionar mi español.
\end{quote}

\begin{quote}
\emph{I want to perfect my Spanish.}
\end{quote}

Built on \emph{la perfección} — perfection.

\begin{cousinweb}
\emph{perfeccionar} is built on \emph{perfección}, a learned borrowing from Latin \emph{perfectio} 'completion', from \emph{perficere} 'to finish', built on \emph{facere} 'to make' (Lewis \& Short).
\end{cousinweb}

\begin{grammarlens}[title={What you've built}]
One more word for telling, explaining and describing at B1.
\end{grammarlens}

\subsection*{Your turn}
\begin{itemize}
\raggedright
  \item \emph{Say it:} \emph{perfeccionar}
  \item \emph{Say it:} \emph{perfeccionar}, once more
  \item \emph{Say it:} a sentence of your own with \emph{perfeccionar}
  \item \emph{Recall:} say the Spanish for the trend, then the Spanish for the progress, then say \emph{perfeccionar} again
\end{itemize}

\begin{tcolorbox}[breakable,colback=teal!4,colframe=teal!35!black,title={Before you move on}]
What does perfeccionar mean? (\textquotedblleft{}To perfect, to polish\textquotedblright{}.) Say it once more.
\end{tcolorbox}
\section[el desafío]{el desafío — the challenge}

\label{lesson:ES-C609-desafio}



\begin{quote}
Before the new one: say the Spanish for the progress, then the Spanish for to perfect.
\end{quote}

\subsection*{You'll want to know: el desafío}
\textbf{el desafío} — \textquotedblleft{}the challenge\textquotedblright{}.

\begin{quote}
Aprender un idioma es un desafío.
\end{quote}

\begin{quote}
\emph{Learning a language is a challenge.}
\end{quote}

The verb is \emph{desafiar} (\emph{desafío, desafías}).

\begin{cousinweb}
\emph{desafío} is built on \emph{desafiar}, which is \emph{des-} + \emph{afiar} (EtymDB), an old verb for pledging faith, from \emph{fiar}, inherited from spoken Latin \emph{fidare}, from \emph{fidere} 'to trust'.
\end{cousinweb}

\begin{grammarlens}[title={What you've built}]
One more word for telling, explaining and describing at B1.
\end{grammarlens}

\subsection*{Your turn}
\begin{itemize}
\raggedright
  \item \emph{Say it:} \emph{el desafío}
  \item \emph{Say it:} \emph{el desafío}, once more
  \item \emph{Say it:} a sentence of your own with \emph{el desafío}
  \item \emph{Recall:} say the Spanish for the progress, then the Spanish for to perfect, then say \emph{el desafío} again
\end{itemize}

\begin{tcolorbox}[breakable,colback=teal!4,colframe=teal!35!black,title={Before you move on}]
What does el desafío mean? (\textquotedblleft{}The challenge\textquotedblright{}.) Say it once more.
\end{tcolorbox}
\section[afrontar]{afrontar — to face, to confront}

\label{lesson:ES-C609-afrontar}



\begin{quote}
Before the new one: say the Spanish for to perfect, then the Spanish for the challenge.
\end{quote}

\subsection*{You'll want to know: afrontar}
\textbf{afrontar} — \textquotedblleft{}to face, to confront\textquotedblright{}.

\begin{quote}
Tenemos que afrontar el problema.
\end{quote}

\begin{quote}
\emph{We have to face the problem.}
\end{quote}

\emph{Afrontar la realidad} — to face reality.

\begin{cousinweb}
\emph{afrontar} is inherited from spoken Latin \emph{affrontare} 'to meet face to face', built on \emph{frons, frontis} 'forehead, front' (EtymDB).
\end{cousinweb}

\begin{grammarlens}[title={What you've built}]
One more word for telling, explaining and describing at B1.
\end{grammarlens}

\subsection*{Your turn}
\begin{itemize}
\raggedright
  \item \emph{Say it:} \emph{afrontar}
  \item \emph{Say it:} \emph{afrontar}, once more
  \item \emph{Say it:} a sentence of your own with \emph{afrontar}
  \item \emph{Recall:} say the Spanish for to perfect, then the Spanish for the challenge, then say \emph{afrontar} again
\end{itemize}

\begin{tcolorbox}[breakable,colback=teal!4,colframe=teal!35!black,title={Before you move on}]
What does afrontar mean? (\textquotedblleft{}To face, to confront\textquotedblright{}.) Say it once more.
\end{tcolorbox}
\section[perseverar]{perseverar — to persevere, to keep at it}

\label{lesson:ES-C609-perseverar}



\begin{quote}
Before the new one: say the Spanish for the challenge, then the Spanish for to face.
\end{quote}

\subsection*{You'll want to know: perseverar}
\textbf{perseverar} — \textquotedblleft{}to persevere, to keep at it\textquotedblright{}.

\begin{quote}
Si perseveras, lo vas a conseguir.
\end{quote}

\begin{quote}
\emph{If you keep at it, you will manage.}
\end{quote}

\emph{La perseverancia} — perseverance.

\begin{cousinweb}
\emph{perseverar} is a learned borrowing from Latin \emph{perseverare} 'to persist' (EtymDB), usually traced to \emph{severus} 'strict, serious'.
\end{cousinweb}

\begin{grammarlens}[title={What you've built}]
One more word for telling, explaining and describing at B1.
\end{grammarlens}

\subsection*{Your turn}
\begin{itemize}
\raggedright
  \item \emph{Say it:} \emph{perseverar}
  \item \emph{Say it:} \emph{perseverar}, once more
  \item \emph{Say it:} a sentence of your own with \emph{perseverar}
  \item \emph{Recall:} say the Spanish for the challenge, then the Spanish for to face, then say \emph{perseverar} again
\end{itemize}

\begin{tcolorbox}[breakable,colback=teal!4,colframe=teal!35!black,title={Before you move on}]
What does perseverar mean? (\textquotedblleft{}To persevere, to keep at it\textquotedblright{}.) Say it once more.
\end{tcolorbox}
